\chapter{Frequency Response and Bode Plots}
\label{ch:bode}
The linear-systems chapter showed \emph{what} first- and second-order frequency
responses look like. This chapter is about \emph{reading and using} them: how to
sketch a Bode plot from a transfer function's poles and zeros, and how to pull a
circuit's behavior back off the plot.

\section{What a Bode Plot Shows}
A Bode plot is a pair of graphs against logarithmic frequency,
\[
20\log_{10}|G(\jj\omega)|
\quad\text{(magnitude, in dB)},
\qquad
\angle G(\jj\omega)
\quad\text{(phase)}.
\]
The magnitude curve shows gain or attenuation; the phase curve shows lead or lag.
Frequency is logarithmic, so equal ratios take equal space: a \emph{decade} is a
\(\times 10\) ratio and an \emph{octave} is a \(\times 2\) ratio. Slopes are quoted
per decade or per octave, and \(\SI{20}{\decibel}\)/decade
\(\approx\SI{6}{\decibel}\)/octave.

\begin{controlsbox}
\textbf{Refresher: the state-space bridge.} Frequency response is the steady
sinusoidal response of a linear system. In state-space form it is
\[
G(\jj\omega)=C(\jj\omega I-A)^{-1}B+D,
\]
i.e.\ the transfer function of \cref{ch:linsys} evaluated on the imaginary axis.
A conceptual bridge, not an exam calculation.
\end{controlsbox}

\section{Building the Sketch from Poles and Zeros}
A Bode magnitude plot can be sketched by hand from the pole and zero locations,
because each contributes a straight-line asymptote that adds to the others:
\begin{itemize}
\item Start at the low-frequency gain (the DC value \(|G(0)|\)).
\item At each \textbf{pole} corner, bend the slope \emph{down} by
\(\SI{20}{\decibel}\) per decade and subtract up to \(90^\circ\) of phase.
\item At each \textbf{zero} corner, bend the slope \emph{up} by
\(\SI{20}{\decibel}\) per decade and add up to \(90^\circ\) of phase.
\item A corner sits at the frequency equal to that pole's or zero's distance from
the origin; a second-order pole pair changes slope by \(\SI{40}{\decibel}\) per
decade near \(\omega_0\).
\end{itemize}
The exact curve rounds the corners---down \(\SI{3}{\decibel}\) at a first-order
pole (\cref{fig:bode1}), and peaking up by \(\sim Q\) at a lightly damped
second-order pair (\cref{fig:bode2})---but the asymptotes get the shape right.

\section{Reading a Circuit off the Plot}
Given a measured or computed Bode plot, four features tell you what the circuit
does:
\begin{description}
\item[Passband and cutoff] the flat region is the passband; the
\(-\SI{3}{\decibel}\) frequency (magnitude \(1/\sqrt2\)) is the cutoff or corner.
\item[Roll-off slope] the far-out slope gives the order: \(-20n\,\si{\decibel}\)
per decade means \(n\) poles. Steeper skirts reject out-of-band signals harder.
\item[Resonant peak] a bump near \(\omega_0\) signals a high-\(Q\) pole pair; its
height grows with \(Q\) and its \(-\SI{3}{\decibel}\) width is the bandwidth
\(BW\approx f_0/Q\).
\item[Phase] the phase slope near a corner flags how much lead or lag the stage
adds---the quantity that governs stability once the stage sits inside a feedback
loop (\cref{ch:linsys}).
\end{description}

\begin{physicalbox}
\textbf{Using it in practice.} To discuss a receiver's IF filter, for example,
read straight off its Bode magnitude: the flat top is the passband, its width at
\(-\SI{3}{\decibel}\) is the bandwidth, the slope past the edges is how fast
adjacent channels are rejected, and any ripple or peak warns of high-\(Q\) stages
that will ring. Every one of these is a pole-zero statement in disguise.
\end{physicalbox}

\section{Bandwidth and Half-Power Points}
For a resonant response the half-power points are where the power is half its
peak, i.e.\ the magnitude is down \(1/\sqrt2\):
\[
20\log_{10}\!\left(\frac{1}{\sqrt2}\right)=-\SI{3.0103}{\decibel}.
\]
The span between them is the bandwidth; high \(Q\) means weak damping, a sharper
peak, and a narrower bandwidth---poles pressed against the \(\jj\omega\) axis.

\begin{exambox}
Bode plots may not be drawn on the exam, but their vocabulary is everywhere:
cutoff frequency, gain, attenuation, phase shift, bandwidth, filter slope and
order, and resonant peaking. Read each as a statement about poles and zeros.
\end{exambox}
